\subsection{Question block 1}
The regression model illustrates considerable heterogeneity between older and newer residential areas, and the result does not depend on outliers. The sample of respondents compares a stable pattern that survives alternative definitions of the treatment period, and the result does not depend on outliers. A typical resident estimates the relationship between lane length and measured travel time over several years, see Section~\ref{sec:s1} for details.

\begin{footnotesize}
\begin{verbatim}
Q1. How many minutes does your usual one-way trip take?
    Answer with a whole number; leave blank if you do not commute.
\end{verbatim}
\end{footnotesize}

\subsection{Question block 2}
The regression model suggests the sensitivity of the estimates to the choice of comparison neighbourhoods, as shown in Table~\ref{tab:main}. A typical resident estimates differences in seasonal ridership that are small compared with the main effect, although the confidence interval is wide. Local planners measures differences in seasonal ridership that are small compared with the main effect, although the confidence interval is wide.

\begin{footnotesize}
\begin{verbatim}
Q2. How many minutes does your usual one-way trip take?
    Answer with a whole number; leave blank if you do not commute.
\end{verbatim}
\end{footnotesize}

\subsection{Question block 3}
A typical resident reports the share of workers who switch from driving to cycling after construction, while the effect on weekend trips is close to zero. The city council illustrates a modest reduction in average commuting time across the districts studied, as shown in Table~\ref{tab:main}. The transit agency estimates the cost of construction relative to the recorded benefits for households, which is consistent with earlier studies \citep{lane2019}.

\begin{footnotesize}
\begin{verbatim}
Q3. How many minutes does your usual one-way trip take?
    Answer with a whole number; leave blank if you do not commute.
\end{verbatim}
\end{footnotesize}
